\scanpage{132}
\mspage{65}
contre il y en a entre $\Pi_{d,d'}$ et $\Pi_d$.
Précisons ce point, en considérant un
morphisme
\[
(215)
\]
\diagram{\begin{tikzcd}[ampersand replacement=\&]\Delta_r\arrow[r,"\alpha"]\&\Delta_{r'}\end{tikzcd}}
\uuline{tel que} $\alpha(0)=0$
alors $D_{\Delta_{r'}}\longrightarrow D_{\Delta_r}$ est un morphisme
étale fini, donc aussi $D^*_{\Delta_{r'}}\unsure{\text{---}|\longrightarrow}D^*_{\Delta_r}$ induit
au dessus de $D^*_{\Delta_r}$, donc on trouve
un diagramme $2$-cartésien (qui \deleted{\unsure{suit|vient}}
\deleted{(209)} \unsure{aussi} se déduit de (209)
\[
(216)
\]
\diagram{\begin{tikzcd}[ampersand replacement=\&, column sep=large,row sep=large]
\widetilde\Pi_{\Delta_{r'}}\arrow[r]\arrow[d]\&\widetilde\Pi_{\Delta_r}\arrow[d]\\
\Pi_{\Delta_{r'}}\arrow[r]\&\Pi_{\Delta_r}\arrow[ul,phantom,"2\text{-cart.}" description]
\end{tikzcd}}
)
\margin{Ceci montre que les $\widetilde\Pi_{\Delta_r}\longrightarrow\Pi_{\Delta_r}$ et $\widetilde\Pi\unsure{\text{---}|\longrightarrow}\Pi$ sont connus \unsure{quand on} connaît les $\widetilde\Pi_d\longrightarrow\Pi_d$, \unsure{des projections locales|des fibrations locales}.}
\margin{\deleted{et les $\unsure{\Pi_d|P_d}$ \unsure{comme}}
\deleted{\diagram{\begin{tikzcd}[ampersand replacement=\&]\unsure{\Pi_d|P_d}\arrow[r]\&\unsure{G_d|\Pi_d}\end{tikzcd}}}
$D_d^*!$ sur $D_{\unsure{\Delta_r|\Delta}}$
\deleted{\diagram{\begin{tikzcd}[ampersand replacement=\&]\unsure{D_d}\arrow[r]\&\unsure{D^*_{\Delta_r}}\end{tikzcd}}}}
où les flèches horizontales induisent sur
les $\pi_1$ des applications injectives à
image d'indice fini (on peut dire que
$\Pi_{\Delta_{r'}}\unsure{\text{---}|\longrightarrow}\Pi_{\Delta_r}$ fait de $\Pi_{\Delta_{r'}}$ un
rev. fini de $\Pi_{\Delta_r}$).

Toutes ces relations sont encore de
nature essentiellement formelle et triviale,
et marcheraient dans un site quelconque,
sans l'hypothèse spécifique des croisements
normaux, et sans que le remplacement des
\margin{131}
